\documentclass[10pt,a4paper]{amsart}
\usepackage{amsmath,amssymb,amsthm,graphicx}
% For commutative diagrams we use Xy-pic
\usepackage[all]{xy}

% We use 2cell for 2-commutative diagrams.
\xyoption{2cell}
% Theorem environments.
%
\theoremstyle{plain}
\newtheorem{theorem}[subsection]{Theorem}
\newtheorem{proposition}[subsection]{Proposition}
\newtheorem{lemma}[subsection]{Lemma}

\theoremstyle{definition}
\newtheorem{definition}[subsection]{Definition}
\newtheorem{example}[subsection]{Example}
\newtheorem{exercise}[subsection]{Exercise}
\newtheorem{situation}[subsection]{Situation}

\theoremstyle{remark}
\newtheorem{remark}[subsection]{Remark}
\newtheorem{remarks}[subsection]{Remarks}

\numberwithin{equation}{subsection}

% Macros
%
\def\lim{\mathop{\mathrm{lim}}\nolimits}
\def\colim{\mathop{\mathrm{colim}}\nolimits}
\def\Spec{\mathop{\mathrm{Spec}}}
\def\Hom{\mathop{\mathrm{Hom}}\nolimits}
\def\Ext{\mathop{\mathrm{Ext}}\nolimits}
\def\SheafHom{\mathop{\mathcal{H}\!\mathit{om}}\nolimits}
\def\SheafExt{\mathop{\mathcal{E}\!\mathit{xt}}\nolimits}
\def\Sch{\mathit{Sch}}
\def\Mor{\mathop{\mathrm{Mor}}\nolimits}
\def\Ob{\mathop{\mathrm{Ob}}\nolimits}
\def\Sh{\mathop{\mathit{Sh}}\nolimits}
\def\NL{\mathop{N\!L}\nolimits}
\def\CH{\mathop{\mathrm{CH}}\nolimits}
\def\proetale{{pro\text{-}\acute{e}tale}}
\def\etale{{\acute{e}tale}}
\def\QCoh{\mathit{QCoh}}
\def\Ker{\mathop{\mathrm{Ker}}}
\def\Im{\mathop{\mathrm{Im}}}
\def\Coker{\mathop{\mathrm{Coker}}}
\def\Coim{\mathop{\mathrm{Coim}}}

% Boxtimes
%
\DeclareMathSymbol{\boxtimes}{\mathbin}{AMSa}{"02}

%
% Macros for moduli stacks/spaces
%
\def\QCohstack{\mathcal{QC}\!\mathit{oh}}
\def\Cohstack{\mathcal{C}\!\mathit{oh}}
\def\Spacesstack{\mathcal{S}\!\mathit{paces}}
\def\Quotfunctor{\mathrm{Quot}}
\def\Hilbfunctor{\mathrm{Hilb}}
\def\Curvesstack{\mathcal{C}\!\mathit{urves}}
\def\Polarizedstack{\mathcal{P}\!\mathit{olarized}}
\def\Complexesstack{\mathcal{C}\!\mathit{omplexes}}
% \Pic is the operator that assigns to X its picard group, usage \Pic(X)
% \Picardstack_{X/B} denotes the Picard stack of X over B
% \Picardfunctor_{X/B} denotes the Picard functor of X over B
\def\Pic{\mathop{\mathrm{Pic}}\nolimits}
\def\Picardstack{\mathcal{P}\!\mathit{ic}}
\def\Picardfunctor{\mathrm{Pic}}
\def\Deformationcategory{\mathcal{D}\!\mathit{ef}}
\title{A locally constant sheaf of modules on a connected scheme is constant after an etale covering}
\author{The Stacks Project Authors}
\date{stacks-project:etale-cohomology.tex:local-snapshot-20261011; GNU FDL 1.2 or later}
\begin{document}
\maketitle
\noindent\textit{Editorial note.} All delivered mathematical text is the author's own, copied from the pinned Stacks Project chapter file etale-cohomology.tex; no statement or proof step was rewritten, reordered or removed.\par\medskip
\noindent\textit{Editorial note.} The proof is the author's complete argument. Its only outside input is the definition of locally constant, namely triviality on the members of an etale covering; the author states this at the start of the proof and then carries out the argument.\par\medskip
\noindent\textit{Editorial note.} The label lemma-connected-locally-constant occurs in no other staged Stacks chapter, so no cross-chapter restatement had to be weighed for this card.\par\medskip
\noindent\textit{Editorial note.} The preamble is the chapter's own native preamble excerpt already used for this source: the Xy-pic setup (native preamble lines 14-18) and the native theorem declarations with the mathematical macros (native preamble lines 154-217).\par\medskip
\noindent\textit{Allowed standard prerequisites (editorial).}\begin{itemize}
\item Locally constant sheaf of Lambda-modules on the small etale site. Same chapter, the definition of a locally constant sheaf of Lambda-modules on X\_etale (constant on the members of an etale covering). Supplies the covering \{U\_i -> X\} on whose members the sheaf is constant, which is the starting datum of the counted proof.
\item Constant sheaf underline M associated to a Lambda-module and its restriction to etale open subschemes. Same chapter, the constant sheaf underline M on X\_etale and the identification of its restriction to U\_i with underline M\_i. Supplies the comparison sheaves underline M\_i and the target sheaf underline M of the conclusion.
\item Etale morphisms are open. Same chapter (also used throughout the chapter), the openness of etale morphisms. Makes U\_M = union of the images Im(U\_i -> X) an open subset of X, so that X = coproduct U\_M is an open covering.
\end{itemize}\medskip
\begin{lemma}
\label{lemma-connected-locally-constant}
Let $X$ be a connected scheme. Let $\Lambda$ be a ring and let
$\mathcal{F}$ be a locally constant sheaf of $\Lambda$-modules.
Then there exists a $\Lambda$-module $M$ and an \'etale covering
$\{U_i \to X\}$ such that $\mathcal{F}|_{U_i} \cong \underline{M}|_{U_i}$.
\end{lemma}

\begin{proof}
Choose an \'etale covering
$\{U_i \to X\}$ such that $\mathcal{F}|_{U_i}$ is constant, say
$\mathcal{F}|_{U_i} \cong \underline{M_i}_{U_i}$.
Observe that $U_i \times_X U_j$ is empty if $M_i$ is not isomorphic
to $M_j$.
For each $\Lambda$-module $M$ let $I_M = \{i \in I \mid M_i \cong M\}$.
As \'etale morphisms are open we see that
$U_M = \bigcup_{i \in I_M} \Im(U_i \to X)$
is an open subset of $X$. Then $X = \coprod U_M$ is a disjoint
open covering of $X$. As $X$ is connected only one $U_M$ is nonempty
and the lemma follows.
\end{proof}

\end{document}
